% Options for packages loaded elsewhere
\PassOptionsToPackage{unicode}{hyperref}
\PassOptionsToPackage{hyphens}{url}
\documentclass[
]{article}
\usepackage{xcolor}
\usepackage[margin=1in]{geometry}
\usepackage{amsmath,amssymb}
\setcounter{secnumdepth}{-\maxdimen} % remove section numbering
\usepackage{iftex}
\ifPDFTeX
  \usepackage[T1]{fontenc}
  \usepackage[utf8]{inputenc}
  \usepackage{textcomp} % provide euro and other symbols
\else % if luatex or xetex
  \usepackage{unicode-math} % this also loads fontspec
  \defaultfontfeatures{Scale=MatchLowercase}
  \defaultfontfeatures[\rmfamily]{Ligatures=TeX,Scale=1}
\fi
\usepackage{lmodern}
\ifPDFTeX\else
  % xetex/luatex font selection
\fi
% Use upquote if available, for straight quotes in verbatim environments
\IfFileExists{upquote.sty}{\usepackage{upquote}}{}
\IfFileExists{microtype.sty}{% use microtype if available
  \usepackage[]{microtype}
  \UseMicrotypeSet[protrusion]{basicmath} % disable protrusion for tt fonts
}{}
\makeatletter
\@ifundefined{KOMAClassName}{% if non-KOMA class
  \IfFileExists{parskip.sty}{%
    \usepackage{parskip}
  }{% else
    \setlength{\parindent}{0pt}
    \setlength{\parskip}{6pt plus 2pt minus 1pt}}
}{% if KOMA class
  \KOMAoptions{parskip=half}}
\makeatother
\usepackage{color}
\usepackage{fancyvrb}
\newcommand{\VerbBar}{|}
\newcommand{\VERB}{\Verb[commandchars=\\\{\}]}
\DefineVerbatimEnvironment{Highlighting}{Verbatim}{commandchars=\\\{\}}
% Add ',fontsize=\small' for more characters per line
\usepackage{framed}
\definecolor{shadecolor}{RGB}{248,248,248}
\newenvironment{Shaded}{\begin{snugshade}}{\end{snugshade}}
\newcommand{\AlertTok}[1]{\textcolor[rgb]{0.94,0.16,0.16}{#1}}
\newcommand{\AnnotationTok}[1]{\textcolor[rgb]{0.56,0.35,0.01}{\textbf{\textit{#1}}}}
\newcommand{\AttributeTok}[1]{\textcolor[rgb]{0.13,0.29,0.53}{#1}}
\newcommand{\BaseNTok}[1]{\textcolor[rgb]{0.00,0.00,0.81}{#1}}
\newcommand{\BuiltInTok}[1]{#1}
\newcommand{\CharTok}[1]{\textcolor[rgb]{0.31,0.60,0.02}{#1}}
\newcommand{\CommentTok}[1]{\textcolor[rgb]{0.56,0.35,0.01}{\textit{#1}}}
\newcommand{\CommentVarTok}[1]{\textcolor[rgb]{0.56,0.35,0.01}{\textbf{\textit{#1}}}}
\newcommand{\ConstantTok}[1]{\textcolor[rgb]{0.56,0.35,0.01}{#1}}
\newcommand{\ControlFlowTok}[1]{\textcolor[rgb]{0.13,0.29,0.53}{\textbf{#1}}}
\newcommand{\DataTypeTok}[1]{\textcolor[rgb]{0.13,0.29,0.53}{#1}}
\newcommand{\DecValTok}[1]{\textcolor[rgb]{0.00,0.00,0.81}{#1}}
\newcommand{\DocumentationTok}[1]{\textcolor[rgb]{0.56,0.35,0.01}{\textbf{\textit{#1}}}}
\newcommand{\ErrorTok}[1]{\textcolor[rgb]{0.64,0.00,0.00}{\textbf{#1}}}
\newcommand{\ExtensionTok}[1]{#1}
\newcommand{\FloatTok}[1]{\textcolor[rgb]{0.00,0.00,0.81}{#1}}
\newcommand{\FunctionTok}[1]{\textcolor[rgb]{0.13,0.29,0.53}{\textbf{#1}}}
\newcommand{\ImportTok}[1]{#1}
\newcommand{\InformationTok}[1]{\textcolor[rgb]{0.56,0.35,0.01}{\textbf{\textit{#1}}}}
\newcommand{\KeywordTok}[1]{\textcolor[rgb]{0.13,0.29,0.53}{\textbf{#1}}}
\newcommand{\NormalTok}[1]{#1}
\newcommand{\OperatorTok}[1]{\textcolor[rgb]{0.81,0.36,0.00}{\textbf{#1}}}
\newcommand{\OtherTok}[1]{\textcolor[rgb]{0.56,0.35,0.01}{#1}}
\newcommand{\PreprocessorTok}[1]{\textcolor[rgb]{0.56,0.35,0.01}{\textit{#1}}}
\newcommand{\RegionMarkerTok}[1]{#1}
\newcommand{\SpecialCharTok}[1]{\textcolor[rgb]{0.81,0.36,0.00}{\textbf{#1}}}
\newcommand{\SpecialStringTok}[1]{\textcolor[rgb]{0.31,0.60,0.02}{#1}}
\newcommand{\StringTok}[1]{\textcolor[rgb]{0.31,0.60,0.02}{#1}}
\newcommand{\VariableTok}[1]{\textcolor[rgb]{0.00,0.00,0.00}{#1}}
\newcommand{\VerbatimStringTok}[1]{\textcolor[rgb]{0.31,0.60,0.02}{#1}}
\newcommand{\WarningTok}[1]{\textcolor[rgb]{0.56,0.35,0.01}{\textbf{\textit{#1}}}}
\usepackage{graphicx}
\makeatletter
\newsavebox\pandoc@box
\newcommand*\pandocbounded[1]{% scales image to fit in text height/width
  \sbox\pandoc@box{#1}%
  \Gscale@div\@tempa{\textheight}{\dimexpr\ht\pandoc@box+\dp\pandoc@box\relax}%
  \Gscale@div\@tempb{\linewidth}{\wd\pandoc@box}%
  \ifdim\@tempb\p@<\@tempa\p@\let\@tempa\@tempb\fi% select the smaller of both
  \ifdim\@tempa\p@<\p@\scalebox{\@tempa}{\usebox\pandoc@box}%
  \else\usebox{\pandoc@box}%
  \fi%
}
% Set default figure placement to htbp
\def\fps@figure{htbp}
\makeatother
\setlength{\emergencystretch}{3em} % prevent overfull lines
\providecommand{\tightlist}{%
  \setlength{\itemsep}{0pt}\setlength{\parskip}{0pt}}
\usepackage{bookmark}
\IfFileExists{xurl.sty}{\usepackage{xurl}}{} % add URL line breaks if available
\urlstyle{same}
\hypersetup{
  hidelinks,
  pdfcreator={LaTeX via pandoc}}

\author{}
\date{\vspace{-2.5em}}

\begin{document}

\subsection{Lab: Food Webs \& Network Measures (Student
Version)}\label{lab-food-webs-network-measures-student-version}

\subsubsection{Lesson Overview}\label{lesson-overview}

\textbf{Computational Topics} - Build and visualize food webs - Write
functions to implement mathematical equations

\textbf{Conservation topics} - Paleofood webs - Species extinction

In this lab we will practice our network visualization and manipulation
skills using the paleo food web data from
\href{https://doi.org/10.1073/pnas.1408471111}{Yeakel et al.~2014}.

\begin{figure}
\centering
\pandocbounded{\includegraphics[keepaspectratio]{labs/figures/paleoweb.jpg}}
\caption{Paleoweb}
\end{figure}

See the beautiful, animated version of the graphic above
\href{https://infograficos.estadao.com.br/public/cidades/extincoes-egito/}{here}

With some interaction networks we can observe the interactions, for
example plant-pollinator networks, seed-disperal networks, human social
networks. In food webs sometimes feeding interactions are observed
directly, through camera traps, people doing timed observations, and now
molecular analysis of gut contents/scat. However, often with food webs
people build probabilistic models of who interacts with who based on
body size (as in the Yeakel et al.~2014), especially with paleowebs.
Thus the data from Yeakel et al.~is 1) an occurrence matrix (Figure 2
from the publication) and a matrix of body sizes (two columns, females
then males). We will use these data to build the foodwebs for each time
period. This lab is pretty challenging because it will use many of our
core programming skills (for loops, writing functions, subsetting data)
and our network skills.

First we will read in the data. The matrix we are reading in has no row
or column names, we will have to set them.

\begin{Shaded}
\begin{Highlighting}[]
\NormalTok{sp\_occ }\OtherTok{\textless{}{-}} \FunctionTok{read.table}\NormalTok{(}\AttributeTok{file=}\StringTok{"data/egypt\_data.txt"}\NormalTok{, }\AttributeTok{header =} \ConstantTok{FALSE}\NormalTok{)}
\FunctionTok{str}\NormalTok{(sp\_occ)}
\end{Highlighting}
\end{Shaded}

\begin{verbatim}
## 'data.frame':    39 obs. of  23 variables:
##  $ V1 : int  1 1 1 1 1 1 1 1 0 1 ...
##  $ V2 : int  1 1 1 1 1 1 1 1 0 1 ...
##  $ V3 : int  1 1 1 1 0 1 1 1 0 1 ...
##  $ V4 : int  1 1 1 1 0 1 1 1 0 1 ...
##  $ V5 : int  1 1 1 1 0 1 1 1 0 1 ...
##  $ V6 : int  1 1 1 1 0 1 1 1 0 1 ...
##  $ V7 : int  1 1 1 1 0 1 1 1 0 1 ...
##  $ V8 : int  1 1 1 1 0 0 1 1 0 1 ...
##  $ V9 : int  1 1 1 1 0 0 1 1 0 0 ...
##  $ V10: int  1 1 1 1 0 0 1 1 0 0 ...
##  $ V11: int  1 1 1 1 0 0 1 1 0 0 ...
##  $ V12: int  1 1 1 1 0 0 1 1 1 0 ...
##  $ V13: int  1 1 1 1 0 0 1 1 1 0 ...
##  $ V14: int  1 1 0 1 0 0 1 1 1 0 ...
##  $ V15: int  1 1 0 1 0 0 1 1 1 0 ...
##  $ V16: int  1 1 0 1 0 0 1 1 1 0 ...
##  $ V17: int  1 1 0 1 0 0 1 1 1 0 ...
##  $ V18: int  1 1 0 1 0 0 0 1 1 0 ...
##  $ V19: int  1 1 0 1 0 0 0 1 1 0 ...
##  $ V20: int  1 1 0 1 0 0 0 1 1 0 ...
##  $ V21: int  1 1 0 1 0 0 0 1 0 0 ...
##  $ V22: int  1 1 0 1 0 0 0 1 0 0 ...
##  $ V23: int  1 1 0 1 0 0 0 0 0 0 ...
\end{verbatim}

\begin{Shaded}
\begin{Highlighting}[]
\NormalTok{sp\_mass }\OtherTok{\textless{}{-}} \FunctionTok{read.table}\NormalTok{(}\AttributeTok{file=}\StringTok{"data/egypt\_mass.txt"}\NormalTok{, }\AttributeTok{header=}\ConstantTok{FALSE}\NormalTok{)}
\FunctionTok{str}\NormalTok{(sp\_mass)}
\end{Highlighting}
\end{Shaded}

\begin{verbatim}
## 'data.frame':    39 obs. of  2 variables:
##  $ V1: int  6 4 18 25 40 122 122 50 35 2200 ...
##  $ V2: int  15 8 36 55 90 260 260 60 65 6300 ...
\end{verbatim}

\begin{figure}
\centering
\pandocbounded{\includegraphics[keepaspectratio]{labs/figures/figure2.jpg}}
\caption{Figure 2}
\end{figure}

The rows are arranged in the order of Figure 2 of the manuscript. To set
the rownames we can make a vector of the names then use the function
\texttt{rownames}. We also have to note which species are predators (all
those in the species in the Carnivora clade in figure 2). Otherwise we
will create a web where giraffes are voracious predators consuming all
of the other species (I made this mistake when constructing the networks
originally). I have transcribed the data from figure 2 for you:

\begin{Shaded}
\begin{Highlighting}[]
\NormalTok{row\_labs\_sp }\OtherTok{\textless{}{-}} \FunctionTok{c}\NormalTok{(}\StringTok{"Canis aureus"}\NormalTok{, }\StringTok{"Vulpes vulpes"}\NormalTok{, }\StringTok{"Lycaon pictus"}\NormalTok{, }\StringTok{"Hyaena hyaena"}\NormalTok{, }\StringTok{"Crocuta crocuta"}\NormalTok{, }\StringTok{"Panthera leo (long maned)"}\NormalTok{, }\StringTok{"Panthera leo (short maned)"}\NormalTok{, }\StringTok{"Panthera pardus"}\NormalTok{, }\StringTok{"Acinonyx jubatus"}\NormalTok{, }\StringTok{"Loxodonta africana"}\NormalTok{, }\StringTok{"Equus asinus"}\NormalTok{, }\StringTok{"Equus grevyi"}\NormalTok{, }\StringTok{"Equus quagga"}\NormalTok{, }\StringTok{"Diceros/Ceratotherium"}\NormalTok{, }\StringTok{"Sus scrofa"}\NormalTok{,  }\StringTok{"Phacochoerus aethiopicus"}\NormalTok{, }\StringTok{"Hippopotamus amphibius"}\NormalTok{, }\StringTok{"Giraffa camelopardalis"}\NormalTok{, }\StringTok{"Dama mesopotamica"}\NormalTok{, }\StringTok{"Camelus dromedarius"}\NormalTok{, }\StringTok{"Taurotragus oryx"}\NormalTok{, }\StringTok{"Tragelaphus spekei"}\NormalTok{, }\StringTok{"Addax nasomaculatus"}\NormalTok{, }\StringTok{"Oryx dammah"}\NormalTok{, }\StringTok{"Oryx beisa"}\NormalTok{, }\StringTok{"Hippotragus equinus"}\NormalTok{, }\StringTok{"Kobus kob"}\NormalTok{, }\StringTok{"Kobus megaceros"}\NormalTok{, }\StringTok{"Alcelaphus bucelaphus"}\NormalTok{, }\StringTok{"Connochaetes taurinus"}\NormalTok{, }\StringTok{"Litocranius walleri"}\NormalTok{, }\StringTok{"Ammodorcas clarkei"}\NormalTok{, }\StringTok{"Gazella dorcas"}\NormalTok{, }\StringTok{"Gazella leptoceros"}\NormalTok{, }\StringTok{"Gazella soemmerringii"}\NormalTok{, }\StringTok{"Capra ibex"}\NormalTok{, }\StringTok{"Ammotragus lervia"}\NormalTok{, }\StringTok{"Bos primigenius"}\NormalTok{, }\StringTok{"Syncerus caffer"}\NormalTok{)}

\DocumentationTok{\#\#\#  Set 1 for predators, 0 for prey  }
\NormalTok{carnivores }\OtherTok{\textless{}{-}} \FunctionTok{c}\NormalTok{(}\FunctionTok{rep}\NormalTok{(}\DecValTok{1}\NormalTok{, }\DecValTok{9}\NormalTok{), }\FunctionTok{rep}\NormalTok{(}\DecValTok{0}\NormalTok{, }\FunctionTok{length}\NormalTok{(row\_labs\_sp)}\SpecialCharTok{{-}} \DecValTok{9}\NormalTok{))}
\FunctionTok{names}\NormalTok{(carnivores) }\OtherTok{\textless{}{-}}\NormalTok{ row\_labs\_sp}
\end{Highlighting}
\end{Shaded}

\subsubsection{Lab question 1: Creating our foodwebs based on body
sizes.}\label{lab-question-1-creating-our-foodwebs-based-on-body-sizes.}

\begin{itemize}
\tightlist
\item
  1a. Use the above vector of species names to label the row names of
  the species occurrence and the body size matrices. The columns of the
  species occurrence matrix are time points, so we can leave those as V1
  etc., but we should set the column names of the mass matrix as ``f'',
  ``m'' (female and male). Use \texttt{head} to check each matrix to see
  if the names are displayed properly.
\end{itemize}

\begin{Shaded}
\begin{Highlighting}[]
\DocumentationTok{\#\#\#  species occurance + mass}
\FunctionTok{rownames}\NormalTok{(sp\_occ) }\OtherTok{\textless{}{-}}\NormalTok{ row\_labs\_sp}
\FunctionTok{rownames}\NormalTok{(sp\_mass) }\OtherTok{\textless{}{-}}\NormalTok{ row\_labs\_sp}

\DocumentationTok{\#\#\# species occurance at time points}
\FunctionTok{colnames}\NormalTok{(sp\_mass) }\OtherTok{\textless{}{-}} \FunctionTok{c}\NormalTok{(}\StringTok{"f"}\NormalTok{, }\StringTok{"m"}\NormalTok{)}

\FunctionTok{head}\NormalTok{(sp\_occ)}
\end{Highlighting}
\end{Shaded}

\begin{verbatim}
##                           V1 V2 V3 V4 V5 V6 V7 V8 V9
## Canis aureus               1  1  1  1  1  1  1  1  1
## Vulpes vulpes              1  1  1  1  1  1  1  1  1
## Lycaon pictus              1  1  1  1  1  1  1  1  1
## Hyaena hyaena              1  1  1  1  1  1  1  1  1
## Crocuta crocuta            1  1  0  0  0  0  0  0  0
## Panthera leo (long maned)  1  1  1  1  1  1  1  0  0
##                           V10 V11 V12 V13 V14 V15 V16
## Canis aureus                1   1   1   1   1   1   1
## Vulpes vulpes               1   1   1   1   1   1   1
## Lycaon pictus               1   1   1   1   0   0   0
## Hyaena hyaena               1   1   1   1   1   1   1
## Crocuta crocuta             0   0   0   0   0   0   0
## Panthera leo (long maned)   0   0   0   0   0   0   0
##                           V17 V18 V19 V20 V21 V22 V23
## Canis aureus                1   1   1   1   1   1   1
## Vulpes vulpes               1   1   1   1   1   1   1
## Lycaon pictus               0   0   0   0   0   0   0
## Hyaena hyaena               1   1   1   1   1   1   1
## Crocuta crocuta             0   0   0   0   0   0   0
## Panthera leo (long maned)   0   0   0   0   0   0   0
\end{verbatim}

\begin{Shaded}
\begin{Highlighting}[]
\FunctionTok{head}\NormalTok{(sp\_mass)}
\end{Highlighting}
\end{Shaded}

\begin{verbatim}
##                             f   m
## Canis aureus                6  15
## Vulpes vulpes               4   8
## Lycaon pictus              18  36
## Hyaena hyaena              25  55
## Crocuta crocuta            40  90
## Panthera leo (long maned) 122 260
\end{verbatim}

Yeakel recommended an updated equation to estimate the probability a
predator consumed a prey based on their relative body masses from
\href{https://doi.org/10.1086/653667}{Rohr et al.~2010.}. The
probability of existence of a trophic link between a predator of
body-size \(m_i\) and a prey of body-size \(m_j\) is given by:

\pandocbounded{\includegraphics[keepaspectratio]{labs/figures/feeding_equ.png}}
(P(\(A_{1j}\) = 1) is the probability predator i eats prey j).

\begin{itemize}
\tightlist
\item
  1b. Write a function and call it \texttt{probEat} to implement the
  equation above. Round the probability to two decimal places.
\end{itemize}

Below are the values of alpha, beta, and gamma for the Serengeti. In
addition, you will need a function to compute the inverse logit function
because this equation is for the logit of the probability, so to
calculate the 0-1 probability you will need to take the inverse logit of
the other side of the equation. Also note, \(log^2\) is equivalent to
(log(\(m_i\)/\(m_j\)))\^{}2

\begin{Shaded}
\begin{Highlighting}[]
\NormalTok{alpha }\OtherTok{\textless{}{-}} \FloatTok{2.51}
\NormalTok{beta }\OtherTok{\textless{}{-}} \FloatTok{0.79}
\NormalTok{gamma }\OtherTok{\textless{}{-}} \SpecialCharTok{{-}}\FloatTok{0.37}
  
\NormalTok{inv\_logit }\OtherTok{\textless{}{-}} \ControlFlowTok{function}\NormalTok{(x) }\FunctionTok{exp}\NormalTok{(x)}\SpecialCharTok{/}\NormalTok{(}\DecValTok{1}\SpecialCharTok{+}\FunctionTok{exp}\NormalTok{(x))}
\NormalTok{inv\_logit}
\end{Highlighting}
\end{Shaded}

\begin{verbatim}
## function (x) 
## exp(x)/(1 + exp(x))
\end{verbatim}

\begin{Shaded}
\begin{Highlighting}[]
\CommentTok{\# your code here}
\NormalTok{inv\_logit }\OtherTok{\textless{}{-}} \ControlFlowTok{function}\NormalTok{(x) }\FunctionTok{exp}\NormalTok{(x)}\SpecialCharTok{/}\NormalTok{(}\DecValTok{1}\SpecialCharTok{+}\FunctionTok{exp}\NormalTok{(x))}

\NormalTok{probEat }\OtherTok{\textless{}{-}} \ControlFlowTok{function}\NormalTok{(pred\_mass, prey\_mass, alpha, beta, gamma) \{}
  \FunctionTok{round}\NormalTok{(}\FunctionTok{inv\_logit}\NormalTok{(alpha }\SpecialCharTok{+}\NormalTok{ beta }\SpecialCharTok{*} \FunctionTok{log}\NormalTok{(pred\_mass }\SpecialCharTok{/}\NormalTok{ prey\_mass)}\SpecialCharTok{+} 
\NormalTok{                    gamma }\SpecialCharTok{*}\NormalTok{ (}\FunctionTok{log}\NormalTok{(pred\_mass }\SpecialCharTok{/}\NormalTok{ prey\_mass))}\SpecialCharTok{\^{}}\DecValTok{2}\NormalTok{), }\DecValTok{2}\NormalTok{)}
\NormalTok{\}}
\end{Highlighting}
\end{Shaded}

\begin{itemize}
\tightlist
\item
  1c. Now create networks of who eats whom. We will start with adjacency
  matrices. We will assume all of our species are the size of females.
  For this step, don`t worry about predators vs.~prey yet, just
  calculate all of the feeding probabilities based on body sizes.
\end{itemize}

Hint: if you start with a square matrix of all zeros (one row and one
column for each species), you can use a for loop to fill in that matrix
with probabilities calculated from your function above.

\begin{Shaded}
\begin{Highlighting}[]
\CommentTok{\# who eats whom, female species}
\NormalTok{n }\OtherTok{\textless{}{-}} \FunctionTok{nrow}\NormalTok{(sp\_mass)}
\NormalTok{prob\_matrix }\OtherTok{\textless{}{-}} \FunctionTok{matrix}\NormalTok{(}\DecValTok{0}\NormalTok{, n, n)}
\FunctionTok{rownames}\NormalTok{(prob\_matrix) }\OtherTok{\textless{}{-}} \FunctionTok{rownames}\NormalTok{(sp\_mass)}
\FunctionTok{colnames}\NormalTok{(prob\_matrix) }\OtherTok{\textless{}{-}} \FunctionTok{rownames}\NormalTok{(sp\_mass)}

\ControlFlowTok{for}\NormalTok{(i }\ControlFlowTok{in} \DecValTok{1}\SpecialCharTok{:}\NormalTok{n)\{}
  \ControlFlowTok{for}\NormalTok{(j }\ControlFlowTok{in} \DecValTok{1}\SpecialCharTok{:}\NormalTok{n)\{}
\NormalTok{    prob\_matrix[i,j] }\OtherTok{\textless{}{-}} \FunctionTok{probEat}\NormalTok{(sp\_mass[i,}\DecValTok{1}\NormalTok{], sp\_mass[j,}\DecValTok{1}\NormalTok{], }\FloatTok{2.51}\NormalTok{, }\FloatTok{0.79}\NormalTok{, }\SpecialCharTok{{-}}\FloatTok{0.37}\NormalTok{)}
\NormalTok{  \}}
\NormalTok{\}}
\FunctionTok{head}\NormalTok{(prob\_matrix)}
\end{Highlighting}
\end{Shaded}

\begin{verbatim}
##                           Canis aureus Vulpes vulpes
## Canis aureus                      0.92          0.94
## Vulpes vulpes                     0.89          0.92
## Lycaon pictus                     0.95          0.95
## Hyaena hyaena                     0.95          0.94
## Crocuta crocuta                   0.94          0.91
## Panthera leo (long maned)         0.82          0.71
##                           Lycaon pictus Hyaena hyaena
## Canis aureus                       0.77          0.65
## Vulpes vulpes                      0.62          0.46
## Lycaon pictus                      0.92          0.90
## Hyaena hyaena                      0.94          0.92
## Crocuta crocuta                    0.95          0.94
## Panthera leo (long maned)          0.94          0.94
##                           Crocuta crocuta
## Canis aureus                         0.42
## Vulpes vulpes                        0.22
## Lycaon pictus                        0.84
## Hyaena hyaena                        0.89
## Crocuta crocuta                      0.92
## Panthera leo (long maned)            0.95
##                           Panthera leo (long maned)
## Canis aureus                                   0.04
## Vulpes vulpes                                  0.01
## Lycaon pictus                                  0.41
## Hyaena hyaena                                  0.58
## Crocuta crocuta                                0.76
## Panthera leo (long maned)                      0.92
##                           Panthera leo (short maned)
## Canis aureus                                    0.04
## Vulpes vulpes                                   0.01
## Lycaon pictus                                   0.41
## Hyaena hyaena                                   0.58
## Crocuta crocuta                                 0.76
## Panthera leo (long maned)                       0.92
##                           Panthera pardus
## Canis aureus                         0.30
## Vulpes vulpes                        0.14
## Lycaon pictus                        0.79
## Hyaena hyaena                        0.86
## Crocuta crocuta                      0.91
## Panthera leo (long maned)            0.95
##                           Acinonyx jubatus
## Canis aureus                          0.49
## Vulpes vulpes                         0.28
## Lycaon pictus                         0.86
## Hyaena hyaena                         0.90
## Crocuta crocuta                       0.93
## Panthera leo (long maned)             0.95
##                           Loxodonta africana
## Canis aureus                            0.00
## Vulpes vulpes                           0.00
## Lycaon pictus                           0.00
## Hyaena hyaena                           0.00
## Crocuta crocuta                         0.00
## Panthera leo (long maned)               0.05
##                           Equus asinus Equus grevyi
## Canis aureus                      0.00         0.00
## Vulpes vulpes                     0.00         0.00
## Lycaon pictus                     0.09         0.04
## Hyaena hyaena                     0.19         0.10
## Crocuta crocuta                   0.41         0.28
## Panthera leo (long maned)         0.84         0.78
##                           Equus quagga
## Canis aureus                      0.01
## Vulpes vulpes                     0.00
## Lycaon pictus                     0.23
## Hyaena hyaena                     0.39
## Crocuta crocuta                   0.63
## Panthera leo (long maned)         0.90
##                           Diceros/Ceratotherium
## Canis aureus                               0.00
## Vulpes vulpes                              0.00
## Lycaon pictus                              0.00
## Hyaena hyaena                              0.01
## Crocuta crocuta                            0.06
## Panthera leo (long maned)                  0.50
##                           Sus scrofa
## Canis aureus                    0.57
## Vulpes vulpes                   0.36
## Lycaon pictus                   0.88
## Hyaena hyaena                   0.91
## Crocuta crocuta                 0.94
## Panthera leo (long maned)       0.95
##                           Phacochoerus aethiopicus
## Canis aureus                                  0.36
## Vulpes vulpes                                 0.17
## Lycaon pictus                                 0.81
## Hyaena hyaena                                 0.87
## Crocuta crocuta                               0.92
## Panthera leo (long maned)                     0.95
##                           Hippopotamus amphibius
## Canis aureus                                0.00
## Vulpes vulpes                               0.00
## Lycaon pictus                               0.01
## Hyaena hyaena                               0.04
## Crocuta crocuta                             0.13
## Panthera leo (long maned)                   0.65
##                           Giraffa camelopardalis
## Canis aureus                                0.00
## Vulpes vulpes                               0.00
## Lycaon pictus                               0.02
## Hyaena hyaena                               0.05
## Crocuta crocuta                             0.17
## Panthera leo (long maned)                   0.70
##                           Dama mesopotamica
## Canis aureus                           0.09
## Vulpes vulpes                          0.03
## Lycaon pictus                          0.57
## Hyaena hyaena                          0.71
## Crocuta crocuta                        0.84
## Panthera leo (long maned)              0.94
##                           Camelus dromedarius
## Canis aureus                             0.00
## Vulpes vulpes                            0.00
## Lycaon pictus                            0.03
## Hyaena hyaena                            0.07
## Crocuta crocuta                          0.22
## Panthera leo (long maned)                0.74
##                           Taurotragus oryx
## Canis aureus                          0.00
## Vulpes vulpes                         0.00
## Lycaon pictus                         0.07
## Hyaena hyaena                         0.15
## Crocuta crocuta                       0.36
## Panthera leo (long maned)             0.82
##                           Tragelaphus spekei
## Canis aureus                            0.42
## Vulpes vulpes                           0.22
## Lycaon pictus                           0.84
## Hyaena hyaena                           0.89
## Crocuta crocuta                         0.92
## Panthera leo (long maned)               0.95
##                           Addax nasomaculatus
## Canis aureus                             0.22
## Vulpes vulpes                            0.09
## Lycaon pictus                            0.74
## Hyaena hyaena                            0.82
## Crocuta crocuta                          0.89
## Panthera leo (long maned)                0.95
##                           Oryx dammah Oryx beisa
## Canis aureus                     0.03       0.04
## Vulpes vulpes                    0.01       0.01
## Lycaon pictus                    0.36       0.44
## Hyaena hyaena                    0.53       0.60
## Crocuta crocuta                  0.73       0.78
## Panthera leo (long maned)        0.92       0.93
##                           Hippotragus equinus
## Canis aureus                             0.01
## Vulpes vulpes                            0.00
## Lycaon pictus                            0.14
## Hyaena hyaena                            0.27
## Crocuta crocuta                          0.52
## Panthera leo (long maned)                0.87
##                           Kobus kob Kobus megaceros
## Canis aureus                   0.22            0.22
## Vulpes vulpes                  0.09            0.09
## Lycaon pictus                  0.74            0.74
## Hyaena hyaena                  0.82            0.82
## Crocuta crocuta                0.89            0.89
## Panthera leo (long maned)      0.95            0.95
##                           Alcelaphus bucelaphus
## Canis aureus                               0.04
## Vulpes vulpes                              0.01
## Lycaon pictus                              0.44
## Hyaena hyaena                              0.60
## Crocuta crocuta                            0.78
## Panthera leo (long maned)                  0.93
##                           Connochaetes taurinus
## Canis aureus                               0.03
## Vulpes vulpes                              0.01
## Lycaon pictus                              0.34
## Hyaena hyaena                              0.51
## Crocuta crocuta                            0.72
## Panthera leo (long maned)                  0.92
##                           Litocranius walleri
## Canis aureus                             0.57
## Vulpes vulpes                            0.36
## Lycaon pictus                            0.88
## Hyaena hyaena                            0.91
## Crocuta crocuta                          0.94
## Panthera leo (long maned)                0.95
##                           Ammodorcas clarkei
## Canis aureus                            0.70
## Vulpes vulpes                           0.52
## Lycaon pictus                           0.91
## Hyaena hyaena                           0.93
## Crocuta crocuta                         0.95
## Panthera leo (long maned)               0.94
##                           Gazella dorcas
## Canis aureus                        0.81
## Vulpes vulpes                       0.69
## Lycaon pictus                       0.93
## Hyaena hyaena                       0.94
## Crocuta crocuta                     0.95
## Panthera leo (long maned)           0.93
##                           Gazella leptoceros
## Canis aureus                            0.83
## Vulpes vulpes                           0.72
## Lycaon pictus                           0.94
## Hyaena hyaena                           0.94
## Crocuta crocuta                         0.95
## Panthera leo (long maned)               0.92
##                           Gazella soemmerringii
## Canis aureus                               0.45
## Vulpes vulpes                              0.24
## Lycaon pictus                              0.85
## Hyaena hyaena                              0.89
## Crocuta crocuta                            0.93
## Panthera leo (long maned)                  0.95
##                           Capra ibex
## Canis aureus                    0.30
## Vulpes vulpes                   0.14
## Lycaon pictus                   0.79
## Hyaena hyaena                   0.86
## Crocuta crocuta                 0.91
## Panthera leo (long maned)       0.95
##                           Ammotragus lervia
## Canis aureus                           0.42
## Vulpes vulpes                          0.22
## Lycaon pictus                          0.84
## Hyaena hyaena                          0.89
## Crocuta crocuta                        0.92
## Panthera leo (long maned)              0.95
##                           Bos primigenius
## Canis aureus                         0.00
## Vulpes vulpes                        0.00
## Lycaon pictus                        0.00
## Hyaena hyaena                        0.00
## Crocuta crocuta                      0.02
## Panthera leo (long maned)            0.31
##                           Syncerus caffer
## Canis aureus                         0.00
## Vulpes vulpes                        0.00
## Lycaon pictus                        0.11
## Hyaena hyaena                        0.22
## Crocuta crocuta                      0.46
## Panthera leo (long maned)            0.85
\end{verbatim}

\begin{itemize}
\tightlist
\item
  1d. Now that you have your matrix of potential feeding interactions
  based on body size, use the \texttt{carnivores} vector created above
  to set all of the feeding interactions of herbivores (0s in that
  vector) to zero. In foodwebs the columns are the higher trophic level
  and the rows are the lower. HINT: the function \texttt{sweep} may be
  useful, though there are many approaches to do the needed matrix
  multiplication. Print the row and column sums.
\end{itemize}

\begin{Shaded}
\begin{Highlighting}[]
\NormalTok{carnivores }\OtherTok{\textless{}{-}} \FunctionTok{c}\NormalTok{(}\FunctionTok{rep}\NormalTok{(}\DecValTok{1}\NormalTok{, }\DecValTok{9}\NormalTok{), }\FunctionTok{rep}\NormalTok{(}\DecValTok{0}\NormalTok{, }\FunctionTok{length}\NormalTok{(row\_labs\_sp)}\SpecialCharTok{{-}} \DecValTok{9}\NormalTok{))}
\CommentTok{\# your code here}
\NormalTok{test }\OtherTok{\textless{}{-}} \FunctionTok{sweep}\NormalTok{(prob\_matrix, }\AttributeTok{MARGIN=}\DecValTok{2}\NormalTok{,carnivores, }\StringTok{\textasciigrave{}}\AttributeTok{*}\StringTok{\textasciigrave{}}\NormalTok{)}
\end{Highlighting}
\end{Shaded}

\subsubsection{Lab question 2: Breaking the networks into time
periods}\label{lab-question-2-breaking-the-networks-into-time-periods}

\begin{itemize}
\tightlist
\item
  2a. With our matrix of feeding interaction we can create a web for
  each time period, including only the species that were not extinct in
  the period. Try first just using the second time period (the second
  column of \texttt{sp\_occ}).
\end{itemize}

Use the function \texttt{empty} from the bipartite package to empty the
matrix of rows and columns with no interactions. The number of species
in the second time period is 36 \texttt{sum(sp\_occ{[},2{]})}. Check to
see that the number of rows in your network with probabilities
\textgreater{} 0 is 36.

HINT: You will need to zero out the rows where a species in not present
in that time period and the columns. The function \texttt{sweep} may be
useful again.

\begin{Shaded}
\begin{Highlighting}[]
\DocumentationTok{\#\#\#  your code here}
\NormalTok{time2 }\OtherTok{\textless{}{-}}\NormalTok{ sp\_occ[, }\DecValTok{2}\NormalTok{]}
\FunctionTok{sum}\NormalTok{(time2)}
\end{Highlighting}
\end{Shaded}

\begin{verbatim}
## [1] 36
\end{verbatim}

\begin{Shaded}
\begin{Highlighting}[]
\NormalTok{web\_time2 }\OtherTok{\textless{}{-}} \FunctionTok{sweep}\NormalTok{(test, }\AttributeTok{MARGIN =} \DecValTok{1}\NormalTok{, time2, }\StringTok{\textasciigrave{}}\AttributeTok{*}\StringTok{\textasciigrave{}}\NormalTok{)}
\NormalTok{web\_time2 }\OtherTok{\textless{}{-}} \FunctionTok{sweep}\NormalTok{(web\_time2, }\AttributeTok{MARGIN =} \DecValTok{2}\NormalTok{, time2, }\StringTok{\textasciigrave{}}\AttributeTok{*}\StringTok{\textasciigrave{}}\NormalTok{)}

\FunctionTok{library}\NormalTok{(bipartite)}
\NormalTok{web\_time2 }\OtherTok{\textless{}{-}} \FunctionTok{empty}\NormalTok{(web\_time2)}
\end{Highlighting}
\end{Shaded}

\begin{itemize}
\tightlist
\item
  2b. Now create a network for all of the time points by creating a list
  where each element is a network. You will need to use a for loop, or
  an \texttt{lapply} if you feel like experimenting with apply
  functions. Print the first 5 columns and rows of the 5th time period.
\end{itemize}

HINT: If choosing the for loop route, remember to create an empty list
of a specific length use the function \texttt{vector}. To access a
specific element of a list, use {[}{[}{]}{]}, for example
cool\_list{[}{[}1{]}{]} accesses the first element of the list.

\begin{Shaded}
\begin{Highlighting}[]
\DocumentationTok{\#\#\#  your code here}
\NormalTok{n\_time }\OtherTok{\textless{}{-}} \FunctionTok{ncol}\NormalTok{(sp\_occ)}
\NormalTok{web3 }\OtherTok{\textless{}{-}} \FunctionTok{vector}\NormalTok{(}\StringTok{"list"}\NormalTok{, n\_time)}

\ControlFlowTok{for}\NormalTok{ (i }\ControlFlowTok{in} \DecValTok{1}\SpecialCharTok{:}\NormalTok{n\_time) \{}
\NormalTok{  time\_i }\OtherTok{\textless{}{-}}\NormalTok{ sp\_occ[, i]                       }\CommentTok{\# }
  
\NormalTok{  web\_i }\OtherTok{\textless{}{-}} \FunctionTok{sweep}\NormalTok{(test, }\AttributeTok{MARGIN =} \DecValTok{1}\NormalTok{, time\_i, }\StringTok{\textasciigrave{}}\AttributeTok{*}\StringTok{\textasciigrave{}}\NormalTok{)  }\CommentTok{\# }
  
\NormalTok{  web\_i }\OtherTok{\textless{}{-}} \FunctionTok{sweep}\NormalTok{(web\_i, }\AttributeTok{MARGIN =} \DecValTok{2}\NormalTok{, time\_i, }\StringTok{\textasciigrave{}}\AttributeTok{*}\StringTok{\textasciigrave{}}\NormalTok{)  }\CommentTok{\# }

\NormalTok{web\_i }\OtherTok{\textless{}{-}} \FunctionTok{empty}\NormalTok{(web\_i)                          }\CommentTok{\# }

\NormalTok{  web3[[i]] }\OtherTok{\textless{}{-}}\NormalTok{ web\_i                             }\CommentTok{\# }
\NormalTok{\}}

\CommentTok{\# Show first 5 rows and columns of time period 5}
\NormalTok{web3[[}\DecValTok{5}\NormalTok{]][}\DecValTok{1}\SpecialCharTok{:}\DecValTok{5}\NormalTok{, }\DecValTok{1}\SpecialCharTok{:}\DecValTok{5}\NormalTok{]}
\end{Highlighting}
\end{Shaded}

\begin{verbatim}
##                           Canis aureus Vulpes vulpes
## Canis aureus                      0.92          0.94
## Vulpes vulpes                     0.89          0.92
## Lycaon pictus                     0.95          0.95
## Hyaena hyaena                     0.95          0.94
## Panthera leo (long maned)         0.82          0.71
##                           Lycaon pictus Hyaena hyaena
## Canis aureus                       0.77          0.65
## Vulpes vulpes                      0.62          0.46
## Lycaon pictus                      0.92          0.90
## Hyaena hyaena                      0.94          0.92
## Panthera leo (long maned)          0.94          0.94
##                           Panthera leo (long maned)
## Canis aureus                                   0.04
## Vulpes vulpes                                  0.01
## Lycaon pictus                                  0.41
## Hyaena hyaena                                  0.58
## Panthera leo (long maned)                      0.92
\end{verbatim}

\subsubsection{Lab question 3: Visualize the
networks}\label{lab-question-3-visualize-the-networks}

\begin{itemize}
\tightlist
\item
  3a. Convert the adjacency matrices to igraph class objects using the
  function \texttt{graph\_from\_adjacency\_matrix}. You can use a for
  loop or an lapply. Because these are food webs, set the argument mode
  to ``directed'' and the argument diag to FALSE (this means a species
  cannot consumer members of its own species, i.e., no
  canabalism/self-loops). Also remember that these interactions are
  weighted.
\end{itemize}

\begin{Shaded}
\begin{Highlighting}[]
\FunctionTok{library}\NormalTok{(igraph)}

\NormalTok{common\_species }\OtherTok{\textless{}{-}} \FunctionTok{intersect}\NormalTok{(}\FunctionTok{rownames}\NormalTok{(web\_time2), }\FunctionTok{colnames}\NormalTok{(web\_time2))}
\NormalTok{web\_square }\OtherTok{\textless{}{-}}\NormalTok{ web\_time2[common\_species, common\_species, drop }\OtherTok{=} \ConstantTok{FALSE}\NormalTok{]}

\NormalTok{igraph\_list }\OtherTok{\textless{}{-}} \FunctionTok{graph\_from\_adjacency\_matrix}\NormalTok{(}
\NormalTok{  web\_square,}
   \AttributeTok{mode =} \StringTok{"directed"}\NormalTok{,}
      \AttributeTok{diag =} \ConstantTok{FALSE}\NormalTok{,}
        \AttributeTok{weighted =} \ConstantTok{TRUE}\NormalTok{)}
\end{Highlighting}
\end{Shaded}

\begin{itemize}
\tightlist
\item
  3b. Plot three networks of your choice, using different colors for the
  predators and prey.
\end{itemize}

\begin{Shaded}
\begin{Highlighting}[]
\DocumentationTok{\#\#\#  assign groups as carnivore or herbivore}
\NormalTok{carnivores }\OtherTok{\textless{}{-}} \FunctionTok{c}\NormalTok{(}\FunctionTok{rep}\NormalTok{(}\DecValTok{1}\NormalTok{, }\DecValTok{9}\NormalTok{), }\FunctionTok{rep}\NormalTok{(}\DecValTok{0}\NormalTok{, }\FunctionTok{length}\NormalTok{(row\_labs\_sp) }\SpecialCharTok{{-}} \DecValTok{9}\NormalTok{))}

\FunctionTok{names}\NormalTok{(carnivores) }\OtherTok{\textless{}{-}}\NormalTok{ row\_labs\_sp}

\NormalTok{chosen\_web }\OtherTok{\textless{}{-}} \FunctionTok{c}\NormalTok{(}\DecValTok{7}\NormalTok{, }\DecValTok{11}\NormalTok{, }\DecValTok{23}\NormalTok{)}

\ControlFlowTok{for}\NormalTok{ (i }\ControlFlowTok{in}\NormalTok{ chosen\_web) \{}
\NormalTok{  web\_matrix }\OtherTok{\textless{}{-}}\NormalTok{ web3[[i]]}


\DocumentationTok{\#\# square matrix}

\NormalTok{species }\OtherTok{\textless{}{-}} \FunctionTok{intersect}\NormalTok{(}\FunctionTok{rownames}\NormalTok{(web\_matrix), }\FunctionTok{colnames}\NormalTok{(web\_matrix))}
\NormalTok{web\_matrix }\OtherTok{\textless{}{-}}\NormalTok{ web\_matrix[species,species]}

  
  \DocumentationTok{\#\#\#  convert to a network 3d object}
\NormalTok{g }\OtherTok{\textless{}{-}} \FunctionTok{graph\_from\_adjacency\_matrix}\NormalTok{(web\_matrix, }\AttributeTok{mode =} \StringTok{"directed"}\NormalTok{, }\AttributeTok{weighted =} \ConstantTok{TRUE}\NormalTok{, }\AttributeTok{diag =} \ConstantTok{FALSE}\NormalTok{)}

\NormalTok{col }\OtherTok{\textless{}{-}} \FunctionTok{ifelse}\NormalTok{(carnivores[species] }\SpecialCharTok{==} \DecValTok{1}\NormalTok{, }\StringTok{"red"}\NormalTok{, }\StringTok{"blue"}\NormalTok{)}
  \DocumentationTok{\#\#\#  plot the network}
  \FunctionTok{plot}\NormalTok{(g, }\AttributeTok{vertex.color =}\NormalTok{ col, }\AttributeTok{main =} \FunctionTok{paste}\NormalTok{(}\StringTok{"Food Web {-} Time"}\NormalTok{, i))}
\NormalTok{\}}
\end{Highlighting}
\end{Shaded}

\pandocbounded{\includegraphics[keepaspectratio]{06-lab-foodwebs_student_files/figure-latex/unnamed-chunk-11-1.pdf}}
\pandocbounded{\includegraphics[keepaspectratio]{06-lab-foodwebs_student_files/figure-latex/unnamed-chunk-11-2.pdf}}
\pandocbounded{\includegraphics[keepaspectratio]{06-lab-foodwebs_student_files/figure-latex/unnamed-chunk-11-3.pdf}}

\begin{Shaded}
\begin{Highlighting}[]
\CommentTok{\# your code here}
\DocumentationTok{\#\#\#  assign groups as carnivore or herbivore}
\NormalTok{carnivores }\OtherTok{\textless{}{-}} \FunctionTok{c}\NormalTok{(}\FunctionTok{rep}\NormalTok{(}\DecValTok{1}\NormalTok{, }\DecValTok{9}\NormalTok{), }\FunctionTok{rep}\NormalTok{(}\DecValTok{0}\NormalTok{, }\FunctionTok{length}\NormalTok{(row\_labs\_sp) }\SpecialCharTok{{-}} \DecValTok{9}\NormalTok{))}

\FunctionTok{names}\NormalTok{(carnivores) }\OtherTok{\textless{}{-}}\NormalTok{ row\_labs\_sp}

\NormalTok{chosen\_web }\OtherTok{\textless{}{-}} \FunctionTok{c}\NormalTok{(}\DecValTok{10}\NormalTok{, }\DecValTok{13}\NormalTok{, }\DecValTok{17}\NormalTok{)}

\ControlFlowTok{for}\NormalTok{ (i }\ControlFlowTok{in}\NormalTok{ chosen\_web) \{}
\NormalTok{  web\_matrix }\OtherTok{\textless{}{-}}\NormalTok{ web3[[i]]}


\DocumentationTok{\#\# square matrix}

\NormalTok{species }\OtherTok{\textless{}{-}} \FunctionTok{intersect}\NormalTok{(}\FunctionTok{rownames}\NormalTok{(web\_matrix), }\FunctionTok{colnames}\NormalTok{(web\_matrix))}
\NormalTok{web\_matrix }\OtherTok{\textless{}{-}}\NormalTok{ web\_matrix[species,species]}

  
  \DocumentationTok{\#\#\#  convert to a network 3d object}
\NormalTok{g }\OtherTok{\textless{}{-}} \FunctionTok{graph\_from\_adjacency\_matrix}\NormalTok{(web\_matrix, }\AttributeTok{mode =} \StringTok{"directed"}\NormalTok{, }\AttributeTok{weighted =} \ConstantTok{TRUE}\NormalTok{, }\AttributeTok{diag =} \ConstantTok{FALSE}\NormalTok{)}

\NormalTok{col }\OtherTok{\textless{}{-}} \FunctionTok{ifelse}\NormalTok{(carnivores[species] }\SpecialCharTok{==} \DecValTok{1}\NormalTok{, }\StringTok{"red"}\NormalTok{, }\StringTok{"blue"}\NormalTok{)}
  \DocumentationTok{\#\#\#  plot the network}
  \FunctionTok{plot}\NormalTok{(g, }\AttributeTok{vertex.color =}\NormalTok{ col, }\AttributeTok{main =} \FunctionTok{paste}\NormalTok{(}\StringTok{"Food Web {-} Time"}\NormalTok{, i))}
\NormalTok{\}}
\end{Highlighting}
\end{Shaded}

\pandocbounded{\includegraphics[keepaspectratio]{06-lab-foodwebs_student_files/figure-latex/unnamed-chunk-12-1.pdf}}
\pandocbounded{\includegraphics[keepaspectratio]{06-lab-foodwebs_student_files/figure-latex/unnamed-chunk-12-2.pdf}}
\pandocbounded{\includegraphics[keepaspectratio]{06-lab-foodwebs_student_files/figure-latex/unnamed-chunk-12-3.pdf}}

\begin{Shaded}
\begin{Highlighting}[]
\CommentTok{\# your code here}
\DocumentationTok{\#\#\#  assign groups as carnivore or herbivore}
\NormalTok{carnivores }\OtherTok{\textless{}{-}} \FunctionTok{c}\NormalTok{(}\FunctionTok{rep}\NormalTok{(}\DecValTok{1}\NormalTok{, }\DecValTok{9}\NormalTok{), }\FunctionTok{rep}\NormalTok{(}\DecValTok{0}\NormalTok{, }\FunctionTok{length}\NormalTok{(row\_labs\_sp) }\SpecialCharTok{{-}} \DecValTok{9}\NormalTok{))}

\FunctionTok{names}\NormalTok{(carnivores) }\OtherTok{\textless{}{-}}\NormalTok{ row\_labs\_sp}

\NormalTok{chosen\_web }\OtherTok{\textless{}{-}} \FunctionTok{c}\NormalTok{(}\DecValTok{1}\NormalTok{, }\DecValTok{3}\NormalTok{, }\DecValTok{6}\NormalTok{)}

\ControlFlowTok{for}\NormalTok{ (i }\ControlFlowTok{in}\NormalTok{ chosen\_web) \{}
\NormalTok{  web\_matrix }\OtherTok{\textless{}{-}}\NormalTok{ web3[[i]]}


\DocumentationTok{\#\# square matrix}

\NormalTok{species }\OtherTok{\textless{}{-}} \FunctionTok{intersect}\NormalTok{(}\FunctionTok{rownames}\NormalTok{(web\_matrix), }\FunctionTok{colnames}\NormalTok{(web\_matrix))}
\NormalTok{web\_matrix }\OtherTok{\textless{}{-}}\NormalTok{ web\_matrix[species,species]}

  
  \DocumentationTok{\#\#\#  convert to a network 3d object}
\NormalTok{g }\OtherTok{\textless{}{-}} \FunctionTok{graph\_from\_adjacency\_matrix}\NormalTok{(web\_matrix, }\AttributeTok{mode =} \StringTok{"directed"}\NormalTok{, }\AttributeTok{weighted =} \ConstantTok{TRUE}\NormalTok{, }\AttributeTok{diag =} \ConstantTok{FALSE}\NormalTok{)}

\NormalTok{col }\OtherTok{\textless{}{-}} \FunctionTok{ifelse}\NormalTok{(carnivores[species] }\SpecialCharTok{==} \DecValTok{1}\NormalTok{, }\StringTok{"red"}\NormalTok{, }\StringTok{"blue"}\NormalTok{)}
  \DocumentationTok{\#\#\#  plot the network}
  \FunctionTok{plot}\NormalTok{(g, }\AttributeTok{vertex.color =}\NormalTok{ col, }\AttributeTok{main =} \FunctionTok{paste}\NormalTok{(}\StringTok{"Food Web {-} Time"}\NormalTok{, i))}
\NormalTok{\}}
\end{Highlighting}
\end{Shaded}

\pandocbounded{\includegraphics[keepaspectratio]{06-lab-foodwebs_student_files/figure-latex/unnamed-chunk-13-1.pdf}}
\pandocbounded{\includegraphics[keepaspectratio]{06-lab-foodwebs_student_files/figure-latex/unnamed-chunk-13-2.pdf}}
\pandocbounded{\includegraphics[keepaspectratio]{06-lab-foodwebs_student_files/figure-latex/unnamed-chunk-13-3.pdf}}

\end{document}
